\section*{Capítulo \ref{ch:b1:spherical-astronomy} --- Astronomia esférica}

\begin{solution}{exo:b1:spherical-astronomy:1}
A \omterm{def:b1:spherical-astronomy:horizontal}{altura angular} do polo é igual à \omterm{def:b1:spherical-astronomy:geographic}{latitude}: \ang{52}. Uma estrela nunca se põe se
$\delta > \ang{90} - \ang{52} = \ang{38}$, e nunca nasce se
$\delta < \ang{-38}$.
\end{solution}

\begin{solution}{exo:b1:spherical-astronomy:2}
$\qty{6.75}{h}\times\ang{15} = \ang{101.25}$. $H = \mathrm{LST} - \alpha =
\qty{10}{h} - \qty{6}{h}\,\qty{45}{min} = \qty{3}{h}\,\qty{15}{min}$, ou seja
\ang{48.75}. O \omterm{def:b1:spherical-astronomy:equatorial}{ângulo horário} é positivo: a estrela está a oeste do meridiano, que ela cruzou há três horas e um quarto siderais.
\end{solution}

\begin{solution}{exo:b1:spherical-astronomy:3}
$h = \ang{90} - \varphi + \delta_\odot$: \ang{73.44} no solstício de junho,
\ang{50} nos equinócios, \ang{26.56} no solstício de dezembro.
\end{solution}

\begin{solution}{exo:b1:spherical-astronomy:4}
1900: divisível por 100, mas não por 400, não é bissexto. 2000: divisível por
400, é bissexto. 2100: não é bissexto. O ano juliano excede o \omterm{def:b1:spherical-astronomy:tropical-year}{ano tropical} em
$365.25 - 365.2422 = 0.0078$ dia, \num{0.78} dia por século; o ano gregoriano
em $0.0003$ dia, \num{0.03} dia por século.
\end{solution}

\begin{solution}{exo:b1:spherical-astronomy:5}
$\cos\Delta = \sin^2\ang{45} + \cos^2\ang{45}\cos\ang{90} = 0.5$, logo
$\Delta = \ang{60}$ e a distância é $6371\times\pi/3 = \qty{6672}{km}$.
Ao longo do paralelo, um \omterm{def:b1:spherical-astronomy:spherical-triangle}{círculo menor} de raio $R\cos\ang{45}$, a distância é
$6371\times\cos\ang{45}\times\pi/2 = \qty{7076}{km}$. O \omterm{def:b1:spherical-astronomy:spherical-triangle}{círculo máximo} é
mais curto por cerca de \qty{404}{km}.
\end{solution}

\begin{solution}{exo:b1:spherical-astronomy:6}
$\sin h = \sin\ang{48}\sin\ang{20} + \cos\ang{48}\cos\ang{20}\cos\ang{30}
= 0.2542 + 0.5446 = 0.7987$, $h = \ang{53.0}$.
$\cos\mathrm{A} = (\sin\ang{20} - \sin\ang{48}\sin\ang{53.0})/(\cos\ang{48}
\cos\ang{53.0}) = (0.3420 - 0.5935)/0.4027 = -0.6246$, $\arccos = \ang{128.7}$;
a estrela está a oeste do meridiano ($H > 0$), logo $\mathrm{A} = \ang{360} -
\ang{128.7} = \ang{231.3}$: para o sudoeste, \ang{53} acima do horizonte.
\end{solution}

\begin{solution}{exo:b1:spherical-astronomy:7}
$\cos H_0 = -\tan\ang{60}\tan\ang{23.44} = -1.7321\times0.4335 = -0.7510$,
$H_0 = \ang{138.67}$, duração do dia $2\times138.67/15 = \qty{18.49}{h}$, cerca de
\qty{18}{h}~\qty{29}{min}, como mostra a curva de \ang{60} N no solstício de junho.
\end{solution}

\begin{solution}{exo:b1:spherical-astronomy:8}
A passagem se adianta \qty{3.93}{min} por dia, logo $30\times3.93 =
\qty{118}{min}$ em 30 dias: a estrela cruza o meridiano por volta das 20 h 02 min. Para cruzar ao
meio-dia ela precisa adiantar-se \qty{10}{h} = \qty{600}{min}, o que leva
$600/3.93 \approx 153$ dias.
\end{solution}

\begin{solution}{exo:b1:spherical-astronomy:9}
(a) Em um equinócio a \omterm{def:b1:spherical-astronomy:horizontal}{altura angular} ao meio-dia é \ang{45}: $I = 1361\times\sin\ang{45}
= \qty{962}{W.m^{-2}}$. (b) No solstício de junho $h = \ang{68.44}$ e
$(a/r)^2 = (149.598/152.100)^2 = 0.9674$:
$I = 1361\times0.9674\times\sin\ang{68.44} = \qty{1224}{W.m^{-2}}$. O Sol mais
alto compensa de sobra a maior distância.
\end{solution}

\begin{solution}{exo:b1:spherical-astronomy:10}
$\sin\delta_\odot = \sin\ang{23.44}\sin\ang{60} = 0.3445$,
$\delta_\odot = \ang{20.15}$. $\tan\alpha_\odot = \cos\ang{23.44}\tan\ang{60}
= 1.5890$, $\alpha_\odot = \ang{57.82}$ (no mesmo quadrante que $\lambda_\odot$).
A \omterm{def:b1:spherical-astronomy:equatorial}{ascensão reta} fica atrás por \ang{2.18}, ou seja, $2.18\times4 \approx
\qty{8.7}{min}$ de tempo. Um Sol que se movesse uniformemente ao longo da \omterm{def:b1:spherical-astronomy:ecliptic}{eclíptica} ainda se moveria de forma irregular em \omterm{def:b1:spherical-astronomy:equatorial}{ascensão reta}, a grandeza que fixa o \omterm{def:b1:spherical-astronomy:equatorial}{ângulo horário} e, portanto, o relógio de sol: atrás do Sol médio entre um equinócio e o solstício seguinte, à frente entre um solstício e o equinócio seguinte. Esse é o termo de obliquidade da \omterm{def:b1:spherical-astronomy:solar-time}{equação do tempo}.
\end{solution}

\begin{solution}{exo:b1:spherical-astronomy:11}
Ao nascer $h = 0$, logo a \cref{prop:b1:spherical-astronomy:transform} dá
$\sin\delta = \cos\varphi\cos\mathrm{A}_0$; no solstício de junho
$\delta = \varepsilon$, logo $\cos\mathrm{A}_0 = \sin\varepsilon/\cos\varphi$.
Para $\mathrm{A}_0 = \ang{45}$: $\cos\varphi = \sin\ang{23.44}/\cos\ang{45} =
0.5626$, $\varphi = \ang{55.8}$ N. Acima de $\ang{90} - \varepsilon =
\ang{66.56}$, $\sin\varepsilon/\cos\varphi > 1$: não há \omterm{def:b1:spherical-astronomy:horizontal}{azimute} de nascer,
porque o Sol não se põe naquele dia (sol da meia-noite).
\end{solution}

\begin{solution}{exo:b1:spherical-astronomy:12}
(a) Um fuso esférico de ângulo $\theta$ é a fração $\theta/2\pi$ da esfera:
$\frac{\theta}{2\pi}\times4\pi R^2 = 2\theta R^2$. (b) Os três pares de fusos
têm área total $2\times(2A + 2B + 2C)R^2$; eles cobrem a esfera uma vez,
exceto o triângulo e o seu oposto, cobertos três vezes cada um, ou seja,
duas vezes mais do que o necessário: $4(A + B + C)R^2 = 4\pi R^2 + 4S$, em
que $S$ é a área do triângulo. Daí $S = R^2(A + B + C - \pi)$. (c) Três ângulos
retos: $S = R^2(3\pi/2 - \pi) = \pi R^2/2$, um oitavo da esfera (um octante).
O triângulo da figura: $89.9 + 64.9 + 67.4 - 180 = \ang{42.2}$, isto é,
\num{0.736} radiano: $S = 0.736\,R^2$, \qty{5.9}{\percent} da esfera.
\end{solution}

\begin{solution}{pb:b1:spherical-astronomy:1}
\textbf{1.} Colatitudes $\ang{90} - (\ang{-24.63}) = \ang{114.63}$ e
$\ang{90} - \ang{51.48} = \ang{38.52}$; diferença de \omterm{def:b1:spherical-astronomy:geographic}{longitude} \ang{70.40}.

\textbf{2.} $\cos\Delta = \cos\ang{114.63}\cos\ang{38.52} +
\sin\ang{114.63}\sin\ang{38.52}\cos\ang{70.40} = -0.3260 + 0.1899 =
-0.1362$.

\textbf{3.} $\Delta = \ang{97.83}$; distância $6371\times1.7074 =
\qty{10878}{km}$, cerca de \qty{10900}{km}.

\textbf{4.}
\[
  \cos\mathrm{Az} = \frac{\cos\ang{38.52} - \cos\ang{114.63}\cos\ang{97.83}}
  {\sin\ang{114.63}\sin\ang{97.83}} = \frac{0.7823 - 0.0568}{0.9006} = 0.8056 ,
\]
$\mathrm{Az} = \ang{36.3}$, para nordeste.

\textbf{5.} $\sin\mathrm{Az} = \sin\ang{70.40}\sin\ang{38.52}/\sin\ang{97.83}
= 0.5922$, o que dá \ang{36.3} ou \ang{143.7}; o cosseno positivo da
questão 4 seleciona \ang{36.3}.

\textbf{6.} Não. Um rumo constante cruza todos os meridianos sob o mesmo ângulo
e se enrola em direção a um polo; o rumo do \omterm{def:b1:spherical-astronomy:spherical-triangle}{círculo máximo} muda continuamente
(é \ang{36.3} só no início). A aeronave voaria um caminho mais longo e chegaria
a outro lugar, a menos que corrigisse o rumo.

\textbf{7.} A \omterm{def:b1:spherical-astronomy:horizontal}{altura angular} do \omterm{def:b1:spherical-astronomy:celestial-sphere}{polo celeste} sul é igual ao módulo da \omterm{def:b1:spherical-astronomy:geographic}{latitude}:
\ang{24.63}.

\textbf{8.} $h = \ang{90} - |\ang{-24.63} - (\ang{-16.72})| = \ang{82.09}$;
$\delta > \varphi$, logo Sirius passa \ang{7.91} ao norte do \omterm{def:b1:spherical-astronomy:celestial-sphere}{zênite}.

\textbf{9.} Uma estrela nunca se põe se $\delta < -(\ang{90} - \ang{24.63}) =
\ang{-65.37}$. Canopus, em \ang{-52.70}, não é circumpolar: ela se põe.

\textbf{10.} $\ang{90} - |\ang{-24.63} + \ang{52.70}| = \ang{61.93}$, ao sul
do \omterm{def:b1:spherical-astronomy:celestial-sphere}{zênite}.

\textbf{11.} $\cos H_0 = -\tan(\ang{-24.63})\tan(\ang{-16.72}) =
-(0.4584)(0.3004) = -0.1377$; $H_0 = \ang{97.92}$. Sirius fica acima do horizonte
$2\times97.92/15 = \qty{13.06}{h}$ de \omterm{def:b1:spherical-astronomy:sidereal-time}{tempo sideral}, cerca de \qty{13}{h}.

\textbf{12.} $\cos\mathrm{A}_0 = \sin(\ang{-16.72})/\cos\ang{24.63} =
-0.3165$: Sirius nasce no \omterm{def:b1:spherical-astronomy:horizontal}{azimute} \ang{108.5}, ao sul do leste.

\textbf{13.} Quando $H = 0$, $\mathrm{LST} = \alpha = \qty{6}{h}~\qty{45}{min}~\qty{9}{s}$.

\textbf{14.} $\qty{86400}{s}\times365.2422/366.2422 = \qty{86164.09}{s} =
\qty{23}{h}~\qty{56}{min}~\qty{4}{s}$; a passagem se adianta
\qty{235.9}{s}, \qty{3}{min}~\qty{56}{s}, a cada noite.

\textbf{15.} Precisamos de $12 + n\times24/365.2422 = 6.7525 + 24$, logo
$n = 18.7525/0.065710 = 285.4$ dias. 285 dias depois de 20 de março é
30 de dezembro: Sirius cruza o meridiano à meia-noite média local na noite
de 30 para 31 de dezembro de 2026.

\textbf{16.} Paranal está \ang{70.40} a oeste de Greenwich: o tempo médio local está
atrasado em $70.40/15 = \qty{4.693}{h}$, \qty{4}{h}~\qty{41.6}{min}.

\textbf{17.} \omterm{def:b1:spherical-astronomy:solar-time}{Tempo solar aparente} = \omterm{def:b1:spherical-astronomy:solar-time}{tempo solar médio} + \omterm{def:b1:spherical-astronomy:solar-time}{equação do tempo} =
$\qty{0}{h} - \qty{3}{min}$: o relógio de sol marca cerca de 23 h 57 min.

\textbf{18.} $7\times\qty{3.93}{min} \approx \qty{27.5}{min}$ mais cedo.

\textbf{19.} $+\ang{23.44}$ no solstício de junho, $-\ang{23.44}$ no
solstício de dezembro.

\textbf{20.} Junho: $\ang{90} - |\ang{-24.63} - \ang{23.44}| = \ang{41.93}$.
Dezembro: $\ang{90} - |\ang{-24.63} + \ang{23.44}| = \ang{88.81}$; como
$\delta_\odot = \ang{-23.44} > \varphi$, o Sol passa \ang{1.19} ao norte do
\omterm{def:b1:spherical-astronomy:celestial-sphere}{zênite}.

\textbf{21.} $I = 1361\times(1/0.9837)^2\times\sin\ang{88.81} = 1361\times
1.0334\times0.9998 = \qty{1406}{W.m^{-2}}$.

\textbf{22.} Em junho, o inverno austral. $\cos H_0 = -\tan(\ang{-24.63})
\tan\ang{23.44} = 0.4584\times0.4335 = 0.1988$: o Sol se põe no \omterm{def:b1:spherical-astronomy:equatorial}{ângulo horário}
$H_0 = \ang{78.53}$.

\textbf{23.} $\cos\mathrm{A}_0 = \sin\ang{23.44}/\cos\ang{24.63} = 0.4376$:
$\mathrm{A}_0 = \ang{64.0}$, ao norte do leste.

\textbf{24.} Duração do dia $2\times78.53/15 = \qty{10.47}{h}$; a noite mais longa
do ano no Paranal dura $24 - 10.47 = \qty{13.53}{h}$, isto é,
\textbf{\qty[text-series-to-math=true]{13}{h}~\qty[text-series-to-math=true]{32}{min}}.
\end{solution}
